\documentclass[11pt]{article}
\usepackage[margin=1in]{geometry}
\usepackage{amsmath, amssymb, amsthm}
\usepackage{mathtools}
\usepackage{booktabs}
\usepackage{array}
\usepackage[expansion=false]{microtype}
\usepackage{hyperref}
\hypersetup{colorlinks=true, linkcolor=blue, citecolor=blue, urlcolor=blue}

\theoremstyle{plain}
\newtheorem{proposition}{Proposition}
\newtheorem{theorem}{Theorem}

\theoremstyle{remark}
\newtheorem*{remark}{Remark}
\newtheorem*{observation}{Observation}

\newcommand{\sln}[1]{\mathfrak{sl}(#1,\mathbb{R})}
\newcommand{\TBL}{T_{B\!-\!L}}
\newcommand{\ad}{\operatorname{ad}}
\newcommand{\Tr}{\operatorname{Tr}}

\title{Pythagorean Arithmetic Structure in the Minimal Pati--Salam
Bracket Algebra}
\author{B.~Wallace\thanks{Independent researcher.}}
\date{April 2026 \\ \textit{Pythagorean Structure} \\
Companion to \textit{Colour from Gravity}}

\begin{document}
\maketitle

\begin{abstract}
We record a structural observation about the dimensionless quantities
that arise from the bracket $[e,\omega]$ in $\sln{4}$ underlying the
minimal Pati--Salam (PS) embedding.  All native ratios produced by
the construction --- charge ratios, adjoint eigenvalues, the
order-three saturation coefficient, and structural multiplicities
--- lie in the multiplicative subgroup $\langle 2, 3 \rangle$ of
$\mathbb{Q}^\times$ generated by the primes 2 and 3.  This is the
arithmetic lattice of Pythagorean tuning.  The lattice property is
specific to the minimal PS embedding: the standard SU(5) and SO(10)
constructions introduce factors of 5 through the 3+2 trace split and
the rank-5 weight space respectively, and therefore do not preserve
this structure.  We test whether the broken-phase Standard Model
mass and mixing spectrum retains any imprint of the lattice and
find no statistically significant preference: the structure is an
ultraviolet algebraic artifact, not a low-energy phenomenological
prediction.  The Koide ratio $\tilde h/h = 2/3$ is the single
$\langle 2, 3 \rangle$ quantity that does survive into the broken
phase, and only because it is a derived algebraic relation
preserved by tree-level matching.  We document the lattice
observation, contrast it with SU(5) and SO(10), report the
empirical negative result, and explicitly disclaim any further
phenomenological consequence.  The note is intended as a
structural record, not a derivation of physics from music or
vice versa.
\end{abstract}

%-----------------------------------------------------------------
\section{Statement}
%-----------------------------------------------------------------

Let $\TBL = \mathrm{diag}(1/3, 1/3, 1/3, -1) \in \sln{4}$ be the
$B-L$ generator of the minimal Pati--Salam embedding, normalized so
that the trace of the generator over the fundamental representation
vanishes.  The dimensionless quantities that arise mechanically from
the bracket structure on $\sln{4}$ with this generator are
enumerated in Table~\ref{tab:ps-ratios}.

\begin{table}[h]
\centering
\small
\begin{tabular}{lll}
\toprule
Quantity & Value & Prime structure \\
\midrule
$\TBL$ quark charge & $1/3$ & $3^{-1}$ \\
$|\TBL|$ lepton charge & $1$ & $1$ \\
Eigenvalue gap of $\TBL$ & $4/3$ & $2^2 \cdot 3^{-1}$ \\
Nontrivial eigenvalue of $\ad_{\TBL}$ & $4/3$ & $2^2 \cdot 3^{-1}$ \\
Coefficient in $\ad_{\TBL}^3 = c \, \ad_{\TBL}$ & $16/9$ & $2^4 \cdot 3^{-2}$ \\
Color multiplicity & $3$ & $3^1$ \\
Dimension of zero eigenspace of $\ad_{\TBL}$ & $9$ & $3^2$ \\
BCH--Jacobi generation count & $3$ & $3^1$ \\
\bottomrule
\end{tabular}
\caption{Native dimensionless quantities of the minimal Pati--Salam
bracket algebra and their prime factorizations.}
\label{tab:ps-ratios}
\end{table}

\begin{observation}
\label{obs:lattice}
Every entry in Table~\ref{tab:ps-ratios} lies in the multiplicative
subgroup $\langle 2, 3 \rangle \subset \mathbb{Q}^\times$ generated
by the primes 2 and 3.  Equivalently, each native ratio is of the
form $2^a \cdot 3^b$ for some $a, b \in \mathbb{Z}$.
\end{observation}

This is mechanical.  The $1/3$ charge arises from $1/N_{\mathrm{color}}$
with $N_{\mathrm{color}} = 3$.  The $4/3$ gap is $1 - (-1/3) = 4/3 =
2^2 \cdot 3^{-1}$.  The $16/9$ saturation coefficient is the square
of the gap.  The multiplicities $3$ and $9$ are dimensions of color
and color-squared structures.  No prime greater than $3$ enters the
construction.

\begin{remark}[Connection to Pythagorean tuning]
Pythagorean tuning is the system in which musical intervals are
defined as the ratios in $\langle 2, 3 \rangle$, with $2$ representing
octave equivalence and $3$ the lowest nontrivial harmonic
\cite{Barker1989}.  Reduced into one octave by powers of two, the
six native PS ratios of Table~\ref{tab:ps-ratios} (above unison)
correspond to the unison ($1$), the major second
$9/8 \leftrightarrow 9$ reduced, the perfect fourth
$4/3 \leftrightarrow 1/3$ or $4/3$ reduced, the perfect fifth
$3/2 \leftrightarrow 3$ reduced, and the Pythagorean minor seventh
$16/9$.

This is not a derivation of music from physics, nor of physics from
music.  It is an identification of two systems that happen to live
on the same arithmetic lattice for different mechanistic reasons.
The PS lattice is generated by $N_{\mathrm{color}} = 3$ and the
charge difference $|q_q - q_\ell| = 4/3$ (factoring as
$2^2 \cdot 3^{-1}$); the Pythagorean lattice is generated by octave
doubling and the third harmonic.
\end{remark}

%-----------------------------------------------------------------
\section{Contrast with SU(5) and SO(10)}
%-----------------------------------------------------------------

The arithmetic lattice property of Observation~\ref{obs:lattice} is
\emph{not} a generic feature of grand-unified embeddings.  We
demonstrate this for the standard alternatives.

\paragraph{SU(5).}  In the Georgi--Glashow embedding, the
hypercharge generator on the fundamental $\mathbf{5}$ is
$Y \propto \mathrm{diag}(-2/3, -2/3, -2/3, +1, +1)$.  The trace
constraint $3 \cdot (-2/3) + 2 \cdot (+1) = 0$ forces the
$3 + 2$ split that defines $\mathrm{SU}(3) \times \mathrm{SU}(2)
\subset \mathrm{SU}(5)$.  The eigenvalue gap is
$1 - (-2/3) = 5/3$, which factors as $5 \cdot 3^{-1}$.  The
order-three saturation coefficient (where it exists) is
$(5/3)^2 = 25/9 \in \langle 3, 5 \rangle$.  The fundamental
representation has dimension $5$.

\paragraph{SO(10).}  The SO(10) construction has rank $5$ and
contains the spinor representation of dimension $16 = 2^4$.  The
rank-$5$ Cartan structure introduces $5$ as a structural prime; the
embedding $\mathrm{SO}(10) \supset \mathrm{SU}(5) \times U(1)$
inherits the $5/3$ gap of SU(5).

\begin{table}[h]
\centering
\small
\begin{tabular}{lll}
\toprule
Construction & Native lattice & Reason for primes \\
\midrule
Minimal Pati--Salam ($\sln{4}$, $\TBL$) & $\langle 2, 3 \rangle$
  & $N_{\mathrm{color}} = 3$, charge gap $= 4/3$ \\
SU(5) (Georgi--Glashow) & $\langle 2, 3, 5 \rangle$
  & $3 + 2$ trace split, fundamental dim $5$ \\
SO(10) (spinor embedding) & $\langle 2, 3, 5 \rangle$
  & rank $5$, inherits SU(5) gap $5/3$ \\
\bottomrule
\end{tabular}
\caption{Native arithmetic lattices of three standard unification
embeddings.  Only the minimal Pati--Salam case lies entirely in
$\langle 2, 3 \rangle$.}
\label{tab:lattice-comparison}
\end{table}

\begin{remark}[Just intonation]
Just intonation extends Pythagorean tuning by including the prime
$5$, giving intervals such as the major third $5/4$ and the major
sixth $5/3$ \cite{Barker1989, Sethares2004}.  The arithmetic lattice
of just intonation is $\langle 2, 3, 5 \rangle$, which is precisely
the lattice into which the SU(5) and SO(10) embeddings fall.  The
parallel between the unification chain
$\mathrm{PS} \subset \mathrm{SU}(5)$ or $\mathrm{SO}(10)$
and the tuning chain $\text{Pythagorean} \subset \text{just intonation}$
is again purely arithmetic; we do not assign physical or musical
significance to it.
\end{remark}

%-----------------------------------------------------------------
\section{What is and is not claimed}
%-----------------------------------------------------------------

For clarity we state explicitly what this note does and does not
claim.

\paragraph{What is claimed.}
\begin{enumerate}\itemsep0pt
  \item The native dimensionless ratios produced by the bracket
    structure on $\sln{4}$ with $\TBL$ are exactly the elements of
    $\langle 2, 3 \rangle$ shown in Table~\ref{tab:ps-ratios}.
  \item This arithmetic lattice coincides with the lattice of
    Pythagorean tuning.
  \item The lattice property is specific to the minimal PS embedding.
    SU(5) and SO(10) introduce the prime $5$ for independent group-
    theoretic reasons (the $3+2$ trace split and the rank-$5$
    Cartan structure respectively).
\end{enumerate}

\paragraph{What is not claimed.}
\begin{enumerate}\itemsep0pt
  \item No statement is made that physics is musical, or that music
    is physical, or that either system derives from the other.
  \item No phenomenological prediction follows from the lattice
    observation.  We tested this directly in
    Section~\ref{sec:empirical} and found that the broken-phase
    Standard Model spectrum does not retain a detectable imprint of
    the $\langle 2, 3 \rangle$ lattice.  The structural observation
    is confined to the unbroken algebraic sector.
  \item No claim is made that one should prefer the minimal
    Pati--Salam embedding over SU(5) or SO(10) on the grounds of
    arithmetic simplicity.  The lattice property is a feature, not
    an argument.
  \item No claim of universality is made.  Other gauge embeddings,
    other Lie-algebra constructions, and other physical systems may
    or may not exhibit similar lattice structure; we have not
    surveyed them.
\end{enumerate}

%-----------------------------------------------------------------
\section{Empirical test: does the broken-phase spectrum preserve the lattice?}
\label{sec:empirical}
%-----------------------------------------------------------------

The arithmetic-lattice property of Observation~\ref{obs:lattice} holds
in the unbroken algebraic sector.  A natural empirical question is
whether the broken-phase mass spectrum and mixing angles preserve any
detectable imprint of this structure.  We run the test here and report
the result.

\paragraph{Method.}
For each measured dimensionless ratio $r$, define the lattice distance
\begin{equation}\label{eq:lattice-distance}
  \varepsilon(r;\, \mathcal{L}) \;=\;
  \min_{(a_1, \dots, a_k) \in \mathbb{Z}^k}
  \bigl|\, \log_2 r \;-\; (a_1 \log_2 p_1 + \dots + a_k \log_2 p_k)\,\bigr|,
\end{equation}
where $\mathcal{L} = \langle p_1, \dots, p_k \rangle$ is the
multiplicative subgroup of $\mathbb{Q}^\times$ generated by the
primes $p_1, \dots, p_k$, and the search is over integers $|a_i| \le N$
for some fixed cap $N$.  This is the minimum log-distance from $r$ to
any lattice point.

For each set of measured ratios, we compute the mean of
$\varepsilon$ and compare against a null distribution: $10^4$
ratios drawn uniformly in log scale over the same range as the
data.  A genuine lattice signal would manifest as a
significantly lower mean lattice distance for the data than for
the null, and as a sharper improvement when the
``correct'' lattice is used.

\paragraph{Data.}
We use PDG 2024 central values: charged-lepton pole masses,
$\overline{\mathrm{MS}}$ light-quark masses at $2~\mathrm{GeV}$,
$\overline{\mathrm{MS}}$ $c$ and $b$ masses at their own scales,
the top pole mass, and the CKM magnitudes from the global fit.
We compute three categories of ratio: intra-generation
(within a single doublet), inter-generation (Yukawa hierarchy),
and CKM-derived.

\paragraph{Result.}
Table~\ref{tab:empirical-result} reports the mean lattice
distance against $\langle 2, 3 \rangle$, $\langle 2, 3, 5 \rangle$,
and $\langle 2, 3, 5, 7 \rangle$ for each ratio category at
$N = 5$, alongside the corresponding null-distribution means.

\begin{table}[h]
\centering
\small
\begin{tabular}{lcccccc}
\toprule
& \multicolumn{2}{c}{$\langle 2, 3 \rangle$}
& \multicolumn{2}{c}{$\langle 2, 3, 5 \rangle$}
& \multicolumn{2}{c}{$\langle 2, 3, 5, 7 \rangle$}\\
\cmidrule(lr){2-3}\cmidrule(lr){4-5}\cmidrule(lr){6-7}
Ratio category & Data & Null
& Data & Null
& Data & Null \\
\midrule
Intra-generation & $0.032$ & $0.040$
& $0.008$ & $0.007$
& $0.001$ & $0.002$ \\
Inter-generation & $0.080$ & $0.065$
& $0.006$ & $0.007$
& $0.001$ & $0.002$ \\
CKM-derived      & $0.035$ & $0.049$
& --- & ---
& --- & --- \\
\bottomrule
\end{tabular}
\caption{Mean lattice distance $\langle \varepsilon \rangle$ in
$\log_2$ units for measured Standard Model ratios versus
null distribution (uniform random in same log range).  All
$z$-scores fall in $[-0.7, +0.3]$; no category shows a
statistically significant signal.}
\label{tab:empirical-result}
\end{table}

The data sits at $0.6$--$1.2 \times$ the null mean across all categories
and all lattices.  The largest absolute $z$-score is $-0.71$ for
CKM ratios on $\langle 2, 3 \rangle$ at $N = 7$, which would be
significant only at the $\sim\!\!0.24$ level (one-tailed).

\begin{proposition}[Negative result: no detectable lattice imprint
in the broken phase]
\label{prop:no-imprint}
The dimensionless mass and mixing ratios of the broken-phase
Standard Model do not exhibit a statistically significant preference
for the multiplicative subgroup $\langle 2, 3 \rangle$ of
$\mathbb{Q}^\times$.  The discriminating diagnostic --- whether
adding the prime 5 to the lattice improves data fit more than it
improves the null --- shows no asymmetry: adding primes reduces
lattice distance equally for the data and the null, consistent with
the data being lattice-generic.
\end{proposition}

\paragraph{Interpretation.}
The negative result is itself the clean conclusion.  The
arithmetic-lattice structure of Observation~\ref{obs:lattice} is
confined to the algebraic sector: it is a property of the
\emph{representation theory} of the bracket $[e, \omega]$ in $\sln{4}$,
not of the broken-phase phenomenology.  The arithmetic-lattice
structure does not survive into the broken-phase observables after
symmetry breaking, renormalization-group evolution, and threshold
matching.  We do not isolate which of these mechanisms is responsible
for the loss of lattice structure; the empirical statement is that
the imprint is not detectable in the infrared.

This places the observation firmly in the category of an
\emph{ultraviolet algebraic artifact}: a correct, mechanical fact
about the construction that does not propagate measurably to
infrared observables.  The structural identification with
Pythagorean tuning remains valid as a statement about the algebraic
sector; it does not extend to a prediction about masses or mixing.

The Koide ratio $\tilde h / h = 2/3$ derived in
\cite{WallaceA} is the one place where a $\langle 2, 3 \rangle$
quantity does survive into the broken-phase Yukawa structure, but
it is a single ratio at a fixed scale; it does not generalize
to a lattice-wide structure across the full mass spectrum.

%-----------------------------------------------------------------
\section*{Acknowledgment}
%-----------------------------------------------------------------

The observation reported here was identified through repeated
adversarial compression: an initial broader claim about ``cosmic
musical structure'' was tested against control distributions of
small rationals, narrowed under the recognition that the relevant
arithmetic constraint is membership in $\langle 2, 3 \rangle$ rather
than statistical coincidence, and finally compressed to the present
structural observation contrasted against SU(5) and SO(10).  The
final form is what survives that process.

%-----------------------------------------------------------------
\begin{thebibliography}{9}

\bibitem{Barker1989}
A.~Barker,
\textit{Greek Musical Writings, Vol.~II: Harmonic and Acoustic Theory},
Cambridge University Press, Cambridge (1989).

\bibitem{Sethares2004}
W.~A.~Sethares,
\textit{Tuning, Timbre, Spectrum, Scale}, 2nd ed.,
Springer-Verlag, London (2004).

\bibitem{WallaceA}
B.~Wallace,
\textit{Colour from Gravity: Pati--Salam Structure from the Palatini
Bracket},
2026.

\end{thebibliography}

\end{document}
